\documentclass[11pt]{article}
\usepackage[margin=1in]{geometry}
\usepackage[T1]{fontenc}
\usepackage{lmodern,amsmath,amssymb,amsthm,mathtools,booktabs,array}
\usepackage[colorlinks=true,allcolors=blue]{hyperref}
\usepackage{microtype}
\newtheorem{theorem}{Theorem}[section]
\newtheorem{proposition}[theorem]{Proposition}
\theoremstyle{definition}
\newtheorem{definition}[theorem]{Definition}
\theoremstyle{remark}
\newtheorem{remark}[theorem]{Remark}
\newcommand{\R}{\mathbb R}
\newcommand{\Z}{\mathbb Z}
\newcommand{\eps}{\varepsilon}
\newcommand{\HH}{\mathbb H}
\newcommand{\OO}{\mathcal O}
\newcommand{\vol}{\operatorname{vol}}
\newcommand{\nuE}{\nu_{\eps}}
\title{Algebraic Infinitesimal Probability\\
\large Geometric Content, Conditioning, and Certified Precision}
\author{Working manuscript accompanying \texttt{hyper-lean}}
\date{18 September 2026}
\begin{document}
\maketitle
\begin{abstract}
We develop normalized probability and finite-observable integration over the
ordered rational-function field $\HH=\R(\eps)$, where $\eps$ is positive and
smaller than every positive real. Exact expressions retain lower-dimensional
geometric contributions, while algebraic $\OO$-notation records optional
precision. We present checked square, cube, disk, and solid-ball dart
contexts, including positive point probabilities, endpoint corrections,
and conditioning across dimensions. The geometric content is a rescaled
classical Wills polynomial, not a newly discovered valuation. The companion
Lean development verifies field arithmetic, normalized contexts, selected
geometric identities, and precision rules. It does not yet construct one
compatible geometric probability on an unrestricted event domain. This
manuscript is a synthesis and formalization draft; novelty and general
geometric extensions remain subjects for further review.
\end{abstract}

\section{Purpose and relation to existing work}
An exact point can receive positive infinitesimal probability without
preventing the whole sample space from having probability one. This idea
is established in non-Archimedean probability and numerosity theory
\cite{nap,numerosity}. Our scope is narrower: finite observable partitions
whose region contents are explicit rational functions. An observable region
may represent a continuum of positions; a finite index set is not a claim
that the geometric board contains finitely many points.

The proposed contribution is a carefully delimited algebraic interface
joining normalization, geometric dimension, exact conditional probability,
and certified remainder bounds. None of these ingredients is claimed to
be new in isolation. In particular, intrinsic volumes and the Wills
polynomial supply the geometric coefficients \cite{wills,steiner}.
A claim that this particular synthesis is original requires a fuller
literature review than the initial comparison reported here.

\section{The exact field and the meaning of algebraic}
Let $\HH=\R(\omega)$ be a rational-function field, ordered by declaring
the indeterminate $\omega$ positive infinite. Equivalently, the sign of
a nonzero rational function is the sign of the ratio of its leading
polynomial coefficients. Define $\eps=\omega^{-1}$. Then
\[
  0<\eps<r\quad(r\in\R,\ r>0),\qquad \eps\omega=1.
\]
All arithmetic below takes place in this exact ordered field. Inverses
of mixed expressions, such as $(1+\eps)^{-1}$, are retained as rational
functions. Finite Laurent polynomials alone would not form a field.

Here ``algebraic'' means finite field operations, polynomial identities,
and order comparisons over specified real geometric data. It does not mean
algebraic over $\mathbb Q$: $\pi$, for example, is an allowed real
coefficient. Nor is $\eps$ algebraic over $\R$; it is an indeterminate.
The formalization uses noncomputable real coefficients, so a claim of
general effective computation with arbitrary real inputs is not made.

\begin{remark}[Scope of the word hyperreal]
This field is a non-Archimedean ordered extension, not a model of full
hyperreal transfer. For example, $\eps$ has no square root in $\R(\eps)$:
the order of a squared rational function at zero is even, whereas that of
$\eps$ is one. Real closure would be a further algebraic extension.
No transfer principle, hyperfinite index set, or infinite summation is
assumed in the present probability constructions.
\end{remark}

\section{Normalized observable contexts}
\begin{definition}
A context consists of a finite index set $I$ and contents $c_i\in\HH$
with $c_i\geq0$ and $Z=\sum_{i\in I}c_i>0$. Define
\[
 \int_{\Omega} f := \frac{\sum_{i\in I} f(i)c_i}{Z},\qquad
 P_\Omega(E):=\int_\Omega\mathbf1_E\quad(E\subseteq I).
\]
The ambient context $\Omega$ includes the partition and its contents.
\end{definition}

\begin{theorem}[Normalization and refinement]
The integral is linear, positive, and sends $1$ to $1$. The probability
is finitely additive, lies between zero and one, and obeys
$P_\Omega(E^c)=1-P_\Omega(E)$. Replacing each region by finitely many
subregions whose contents sum to its content preserves integrals of
functions constant on the original region.
\end{theorem}
\begin{proof}
Expand the finite sums. Positivity uses $c_i\geq0$ and $Z>0$.
Refinement follows by regrouping the subregion contents inside each
original summand. No infinite interchange of sums is involved.
\end{proof}

For $P_\Omega(E)>0$, conditioning is explicitly a change of context:
\[
 \int_{\Omega\mid E}f=
 \frac{\int_\Omega\mathbf1_E f}{P_\Omega(E)}.
\]
Ordinary restriction $\int_\Omega\mathbf1_E f$ does not change the
normalizer. In a product of two contexts the contents multiply;
finite algebra gives Fubini's identity and the usual factorization for
independent coordinate events. Independence is a modeling assumption,
not a consequence of rotation invariance.

\section{Geometry, calibration, and endpoint corrections}
Suppose the real coefficient $F(v)$ assigned to a segment displacement
is rotation invariant, satisfies $F(tv)=tF(v)$ for $t\geq0$, and
$F(1,0)=1$. Then
\[
 F(x,y)=\sqrt{x^2+y^2}.
\]
Indeed, for positive length $\ell$, the coefficients $x/\ell,y/\ell$
define an orthogonal rotation taking $(1,0)$ to $(x/\ell,y/\ell)$.
Scaling proves the claim; the zero case follows by scaling by zero.
This characterizes one coefficient, not every possible probability rule.

Fix planar point content $\eps^2$. If closed segments have the graded
form $\ell\eps+k$, with the same endpoint coefficient $k$, joining two
at one endpoint and applying inclusion--exclusion gives
$k=2k-\eps^2$. Thus
\begin{equation}\label{eq:segments}
 \nuE([a,b])=\ell\eps+\eps^2,\qquad
 \nuE([a,b))=\ell\eps,\qquad
 \nuE((a,b))=\ell\eps-\eps^2.
\end{equation}
The latter two expressions concern positive standard lengths. A
degenerate closed segment is a point. These corrections are needed for
exact finite additivity, even when omitted from a leading-order display.

For a closed convex polygon of area $A$ and perimeter $P$, take
\begin{equation}\label{eq:polygon}
 \nuE(K)=A+\frac P2\eps+\eps^2.
\end{equation}
An internal cut of length $L$ contributes twice to the sum of the two
perimeters. Subtracting the shared closed segment cancels that excess
and one point-level term. This checks the elementary dissection identity;
it is not a formal construction of a valuation on all geometric sets.

\subsection{The Wills polynomial connection}
For a compact convex $K\subseteq\R^d$, let $V_j(K)$ denote its intrinsic
volumes in the standard normalization. The Wills polynomial is
$g_K(z)=\sum_j V_j(K)z^j$ \cite{wills}. Our proposed content is precisely
\begin{equation}\label{eq:wills}
 \nuE(K)=\sum_{j=0}^d V_j(K)\eps^{d-j}
        =\eps^d g_K(\eps^{-1}).
\end{equation}
This is formal polynomial evaluation, not a dilation of a standard set
by a nonstandard scalar. Thus the coefficient polynomial is established
geometry; the normalization and precision interface are our subject.
For general event algebras, a compatible extension and positivity proof
are still necessary. A valuation on convex sets alone is not yet a
probability on all subsets of a board.

\section{Square, cube, and coordinate codimension}
For a closed rectangle, equation~\eqref{eq:polygon} factors as
$\nuE([0,a]\times[0,b])=(a+\eps)(b+\eps)$.
The selected orthogonal-product rule therefore gives
\[
 \nuE([0,1]^d)=(1+\eps)^d.
\]
Set $\eta=\eps/(1+\eps)$; this is a change of gauge, not an equality
between $\eta$ and $\eps$.

\begin{theorem}[Coordinate codimension]
If $r$ coordinates in a closed unit $d$-cube are fixed and the other
$d-r$ range over the entire unit interval, the product context gives
\[
 P(E)=\frac{(1+\eps)^{d-r}\eps^r}{(1+\eps)^d}=\eta^r.
\]
\end{theorem}
\begin{proof}
A free coordinate contributes $1+\eps$; a fixed coordinate contributes
$\eps$. Multiply, then divide by the whole-cube content.
\end{proof}

In particular, the cube, a full coordinate slice, a full axis crossing,
and a point have probabilities $1,\eta,\eta^2,\eta^3$. A general closed
segment of length $L$ in the square has
\[
 P(L)=L\eta+(1-L)\eta^2,
\]
so its diagonal has $\sqrt2\,\eta+(1-\sqrt2)\eta^2$. The corresponding
spatial segment has $L\eta^2+(1-L)\eta^3$. The length coefficient is
geometric; the correction records its endpoints.

These statements differ from a coordinate-grid sampling convention,
where a diagonal and an axis row can contain the same number of samples.
That is a different context, not a contradiction with Euclidean length.

\section{Disk and solid-ball darts}
The standard intrinsic-volume normalization is
\[
 \vol(K+tB_d)=\sum_{j=0}^d\kappa_{d-j}V_j(K)t^{d-j},
 \qquad(\kappa_0,\kappa_1,\kappa_2,\kappa_3)
       =(1,2,\pi,4\pi/3),
\]
as in \cite{steiner}. For round bodies, expanding the ordinary radius
polynomials $\pi(r+t)^2$ and $(4\pi/3)(r+t)^3$ and matching coefficients
gives the selected contents
\[
 \nuE(D_r)=\pi r^2+\pi r\eps+\eps^2,\qquad
 \nuE(B_r)=\frac{4\pi r^3}{3}+2\pi r^2\eps+4r\eps^2+\eps^3.
\]
This step uses given real geometric calibrations and finite polynomial
arithmetic. It does not derive ordinary area or $\pi$ from field axioms.
For radius one, write $Z_D=\nuE(D_1)$ and $Z_B=\nuE(B_1)$.
\begin{center}
\begin{tabular}{lll}
\toprule
Board & Event & Exact probability \\
\midrule
Disk & point & $\eps^2/Z_D$ \\
Disk & diameter & $(2\eps+\eps^2)/Z_D$ \\
Ball & point & $\eps^3/Z_B$ \\
Ball & diameter & $(2\eps^2+\eps^3)/Z_B$ \\
Ball & central disk & $(\pi\eps+\pi\eps^2+\eps^3)/Z_B$ \\
\bottomrule
\end{tabular}
\end{center}

These formulas have concrete positive contexts. For the disk, the
successive disjoint regions are a point, the rest of its diameter, and
the rest of the board, with contents
\[
 \eps^2,\quad 2\eps,\quad \pi+(\pi-2)\eps.
\]
For the ball, the nested point, diameter, and central disk yield the
four disjoint-region contents
\[
 \eps^3,\quad 2\eps^2,\quad
 \pi\eps+(\pi-2)\eps^2,\quad
 \frac{4\pi}{3}+\pi\eps+(4-\pi)\eps^2.
\]
They are positive and sum to $Z_D$ and $Z_B$, respectively, using
$2<\pi<4$. The normalized-context theorem therefore applies. In either
board a point on the chosen diameter has the exact conditional probability
\begin{equation}\label{eq:conditional}
 P(\text{point}\mid\text{diameter})=\frac{\eps}{2+\eps}.
\end{equation}
The denominator of the whole board and its transverse powers of $\eps$
cancel. ``Ball'' here means the solid three-dimensional region; sampling
only a spherical surface requires a different, two-dimensional context.

\section{Algebraic precision and normalized integration}
\begin{definition}
For a fixed positive infinitesimal gauge $g$ and $n\in\Z$, write
\[
 r\in\OO(g^n)\ \Longleftrightarrow\
 \exists M\in\R_{\geq0}: |r|\leq Mg^n.
\]
The statement $x=y+\OO(g^n)$ means $x-y\in\OO(g^n)$.
\end{definition}
The bounding coefficient must be standard real. Equivalently,
$\OO(g^n)=g^n\mathcal F$, where $\mathcal F$ is the subring of elements
bounded by standard reals. There is no limiting variable. Integer orders
also accommodate infinite scales. For $0<g\leq1$, the elementary rules are
\[
 \OO(g^m)\OO(g^n)=\OO(g^{m+n}),\quad
 \OO(g^n)/g^m=\OO(g^{n-m}),\quad
 \OO(g^n)\subseteq\OO(g^m)\ (m\leq n).
\]
Addition uses the weaker of the two orders. Reciprocals preserve an
absolute error order when both reciprocals are bounded; an infinitesimal
denominator requires explicit scale tracking.

\begin{proposition}[Uniform precision under normalization]
If $|f(i)-h(i)|\leq Mg^n$ for every region of a positive normalized
context, using one standard real $M\geq0$, then
\[
 \int_\Omega f=\int_\Omega h+\OO(g^n).
\]
\end{proposition}
\begin{proof}
Integrate the inequalities $-Mg^n\leq f-h\leq Mg^n$. Positivity,
linearity, and $\int_\Omega1=1$ give the required two-sided bound.
\end{proof}

For example, a spatial segment has
$P(L)=L\eta^2+\OO(\eta^3)$, with exact expression
$L\eta^2+(1-L)\eta^3$ retained. A closed planar patch of area $A$ and
perimeter $P$ in the unit cube has
\[
 P(\text{patch})=A\eta+(P/2-2A)\eta^2+(A-P/2+1)\eta^3
               =A\eta+\OO(\eta^2).
\]
An $\OO$ statement by itself forgets the omitted coefficients. More
precision must come from the retained exact expression or additional
information. Cancellation cannot justify dropping a remainder: $(1+g)-1=g$,
not zero. The exact identity
\[
 (1+g)^{-1}=1-g+\frac{g^2}{1+g}=1-g+\OO(g^2)
\]
illustrates finite algebraic expansion without an infinite series.

\section{Points, dots, halos, and the event domain}
A singleton, a radius-$\eps$ dot, and a halo are distinct sets. For
example, displacement $\eps/2$ is in the dot but is not its center;
displacement $2\eps$ lies in the halo but not the dot. Here a halo
contains points at every infinitesimal displacement, not merely standard
real multiples of one chosen $\eps$.

The normalized integral can give a singleton positive content directly.
For any admissible support $E$ with $P_\Omega(E)>0$, the density
$\delta_E=\mathbf1_E/P_\Omega(E)$ integrates to one. Its integral against
$f$ is the conditional average on $E$; exact evaluation at a point
requires singleton support (or an integrand constant on that support).
This distinguishes point evaluation from dot averaging without silently
assigning content to an entire halo.

The geometry is deliberately restricted to standard-coordinate data.
Substituting an infinitesimal length into the open-segment formula would
give, for example,
\[
 (\eps/2)\eps-\eps^2=-\eps^2/2<0.
\]
Consequently that formula cannot define a positive content on all
infinitesimal-coordinate regions. A literal infinitesimal dot needs a
separate compatible event construction. This obstruction is part of the
theory's boundary, not a term to discard in an approximation.

\section{Formalization and remaining mathematical work}
The companion code separates exact rational-function arithmetic from the
historical list backend. The current audited path uses the following modules:
\begin{itemize}
\item \path{Hyper/OrderedRational.lean}: field order and infinitesimal bounds.
\item \path{Hyper/ContextIntegral.lean}: finite normalized integrals,
conditioning, refinements, products, and support densities.
\item \path{Hyper/GeometricContent.lean}: segment coefficient characterization,
endpoint forcing, planar dissection algebra, and square formulas.
\item \path{Hyper/CubicContent.lean}: box products, coordinate codimension,
and positive cube contexts.
\item \path{Hyper/RoundContent.lean}: selected positive disk/ball contexts
and the exact ratios in equation~\eqref{eq:conditional}.
\item \path{Hyper/AlgebraicOrder.lean}: remainder arithmetic, uniform
integral bounds, and selected geometric precision statements.
\item \path{Hyper/AlgebraicSupport.lean}: point, dot, and halo distinctions.
\end{itemize}

\begin{center}
\begin{tabular}{p{0.61\linewidth}p{0.27\linewidth}}
\toprule
Claim & Status in the companion \\
\midrule
Exact field and finite normalized-context laws & Lean-checked \\
Square/cube/disk/ball selected event contexts & Lean-checked \\
Endpoint, product, and dissection identities & Lean-checked \\
Algebraic remainder and uniform integration rules & Lean-checked \\
Global geometric valuation independent of decomposition & Not implemented \\
All curved or infinitesimal-coordinate event sets & Not established \\
Full transfer, countable additivity, LLN or CLT & Not claimed \\
Novelty of the synthesis & Not established \\
\bottomrule
\end{tabular}
\end{center}

The important next mathematical step is to specify one geometric event
algebra and construct its compatible content, proving positivity and
independence of decomposition. Local dissection identities and several
finite contexts are evidence for such a development, but do not replace it.
Uniqueness or minimal-axiom claims would require a separate theorem with
all calibration and domain hypotheses stated.

\section*{Reproducibility and manuscript status}
The implementation is available at
\url{https://github.com/pannous/hyper-lean}; the code revision reviewed for
this draft is \texttt{0cfcc3d}. Run \texttt{bash test.sh} with the pinned
Lean/mathlib setup. The script builds before testing and audits the named
exact-path theorems. Their reported dependencies are
\texttt{propext}, \texttt{Classical.choice}, and \texttt{Quot.sound};
no \texttt{sorryAx} or legacy raw-list equality axiom is accepted on that
path. This is a theorem audit, not a claim that all historical modules are
free of unfinished proofs.

AI tools assisted drafting and proof scripting. This draft requires human
mathematical review and a fuller comparison with prior literature before
a novelty claim or journal submission. It supersedes the probability
claims of the older introductory manuscript; it does not validate that
manuscript's broader calculus and transfer claims.

\begin{thebibliography}{9}
\bibitem{nap} V. Benci, L. Horsten, and S. Wenmackers,
\emph{Non-Archimedean Probability},
\url{https://arxiv.org/abs/1106.1524},
\href{https://doi.org/10.1007/s00032-012-0191-x}{doi:10.1007/s00032-012-0191-x}.
\bibitem{numerosity} V. Benci, E. Bottazzi, and M. Di Nasso,
\emph{Elementary numerosity and measures}, Journal of Logic and Analysis
6 (2014), paper 3, \url{https://arxiv.org/abs/1212.6201}.
\bibitem{wills} M. A. Hern\'andez Cifre and J. Yepes Nicol\'as,
\emph{On the roots of the Wills functional},
\href{https://doi.org/10.1016/j.jmaa.2012.12.065}{doi:10.1016/j.jmaa.2012.12.065}.
\bibitem{steiner} M. B. McCoy and J. A. Tropp,
\emph{From Steiner formulas for cones to concentration of intrinsic volumes},
equation (1.1), \url{https://arxiv.org/abs/1308.5265}.
\end{thebibliography}
\end{document}
